% ---------------------------------------------------------------------------
% Validation of the deterministic MILP on TDSP benchmark instances.
% Suggested placement: last subsection of Section \ref{sec:exp} (experimental
% design), or first subsection of Section \ref{sec:results}.
% Requires no new packages (booktabs already loaded).
% ---------------------------------------------------------------------------

\subsection{Validation of the deterministic model on TDSP benchmark instances}
\label{sec:validation}

Before turning to the stochastic experiments, we verify that the HoS core of
the deterministic MILP of Section~\ref{sec:model} reproduces known optimal
solutions of the pure truck driver scheduling problem. To this end we use the
40 R15-PGLT benchmark instances made available by
\citet{garaix_label-setting_2024}\footnote{\url{https://emse.fr/~garaix/TruckDriverSchedule/Expert_Systems/}},
which comprise fixed customer sequences with 2--14 customers, per-stop service
times (including a loading operation at the origin depot), hard time windows,
and all durations rounded to multiples of 15~minutes. We compare against the
published optima of their \emph{no-night} configuration, which enforces
exactly the daily provisions of Regulation~(EC) No~561/2006 and Directive
2002/15/EC considered in this paper, without the night-work rule.

Since these are diesel instances, the benchmark exercises the HoS machinery of
our model in isolation: we disable the battery ($E_i = 0$,
$\mathcal K = \emptyset$) and reduce all maneuver, queue, and out-of-window
penalty terms to zero, with time windows enforced as hard constraints
($\delta_i = 0$). Five semantic differences between the two formulations are
bridged as follows. First, the reference model allows breaks and rests to
preempt driving and service at arbitrary points; we emulate this on the
15-minute grid by inserting rest-area nodes ($\mathcal L$) every 15~minutes of
driving and by splitting each service operation into 15-minute segments, the
time window applying to the start of the first segment and the whole
operation having to complete by the closing of the window plus its duration.
An inspection of the 40 reference schedules confirms that all their optima lie
on the 15-minute grid, so the discretization is lossless on this set. Second,
a zero-length rest-area node inserted before each customer permits the
arrive--rest--serve pattern present in several reference optima. Third, their
schedules must terminate with a daily rest at the destination (regular, or
reduced against the weekly budget $\bar\rho$), whose duration counts toward
the objective; we add the corresponding constraint
$\rho^{r_1}_{j} + \rho^{r_2}_{j} = 1$ at the final stop. Fourth, the departure
time from the origin is left free (bounded below by the earliest start), as
some instances open the depot window after time zero. Fifth, the working-time
provisions of Directive 2002/15/EC --- at most 6~h of cumulated work without a
break of at least 15~min, and at least 30~min (resp.\ 45~min) of cumulated
breaks in any shift whose working time exceeds 6~h (resp.\ 9~h) --- which our
base model verifies ex post, are here enforced in-model by two additional
accumulators mirroring the consecutive-driving construction
\eqref{eq:cd_prop}--\eqref{eq:cd_ub}.

Table~\ref{tab:validation} summarizes the outcome. Our formulation, solved
with Gurobi to zero optimality gap, reproduces the published optimal objective
exactly on 39 of the 40 instances, with solve times below 64~s (median below
1~s) against instance sizes of up to 180 stops after discretization. On the
single divergent instance (Test~4) our model attains 1410~min against a
published value of 1425~min; a manual verification of our schedule against
the reference formulation confirms it satisfies every constraint of that
model, so the published value appears to be marginally suboptimal on this
instance. We conclude that the break and split-break logic, the reduced-rest
and extended-driving budgets, and the shift-spread constraints
\eqref{eq:h_prop}--\eqref{eq:h_ub} --- which we show to be equivalent to the
requirement that a daily rest be completed within 24~h of the end of the
previous one --- are consistent with an independently published exact model.

\begin{table}[htbp]
\centering
\caption{Validation on the 40 R15-PGLT instances (no-night configuration).}
\label{tab:validation}
\small
\begin{tabular}{lr}
\toprule
Instances solved to proven optimality & 40/40 \\
Exact matches with published optima & 39/40 \\
Divergent instances & 1 (ours 1410 vs.\ 1425, see text) \\
Stops after 15-min discretization & 23--180 \\
Solve time (median / max) & $<1$~s / 64~s \\
\bottomrule
\end{tabular}
\end{table}

% ---------------------------------------------------------------------------
% Optional one-liner for a "code and data availability" note or footnote:
%   The benchmark harness, the per-instance comparison table, and the adapted
%   instances are available with our implementation (pglt.py,
%   solutions/pglt_comparison.csv, data/R15-PGLT/).
% Note: the PGLT instances originate from Pe\~na-Arenas et al.; if you want to
% credit the instance authors separately from garaix_label-setting_2024, add
% the corresponding entry to your .bib and cite it in the first paragraph.
% ---------------------------------------------------------------------------
